% ==============================================================================
% CONFIG: COMMANDS.TEX
% Các lệnh tùy biến hỗ trợ sinh viên thao tác nhanh với hình ảnh, bảng biểu,
% ghi chú học thuật và định dạng văn bản.
% ==============================================================================

% ------------------------------------------------------------------------------
% 1. LỆNH CHÈN HÌNH ĐƠN TIÊU CHUẨN
% Cú pháp: \insertfigure[tùy chọn độ rộng]{đường dẫn}{Chú thích hình}{nhãn label}
% Mặc định độ rộng là 0.8\linewidth nếu không chỉ định.
% Ví dụ: \insertfigure[width=0.75\linewidth]{figures/chapter3/ai_pipeline.png}{Sơ đồ kiến trúc}{fig:pipeline}
% ------------------------------------------------------------------------------
\newcommand{\insertfigure}[4][width=0.8\linewidth]{
    \begin{figure}[htbp]
        \centering
        \includegraphics[#1]{#2}
        \caption{#3}
        \label{#4}
    \end{figure}
}

% ------------------------------------------------------------------------------
% 2. LỆNH CHÈN HAI HÌNH SONG SONG (SUBFIGURES)
% Cú pháp:
% \inserttwofigures{ảnh 1}{caption 1}{label 1}{ảnh 2}{caption 2}{label 2}{caption chung}{label chung}
% ------------------------------------------------------------------------------
\newcommand{\inserttwofigures}[8]{
    \begin{figure}[htbp]
        \centering
        \begin{subfigure}[b]{0.48\textwidth}
            \centering
            \includegraphics[width=\linewidth]{#1}
            \caption{#2}
            \label{#3}
        \end{subfigure}
        \hfill
        \begin{subfigure}[b]{0.48\textwidth}
            \centering
            \includegraphics[width=\linewidth]{#4}
            \caption{#5}
            \label{#6}
        \end{subfigure}
        \caption{#7}
        \label{#8}
    \end{figure}
}

% ------------------------------------------------------------------------------
% 3. HỘP GHI CHÚ HỌC THUẬT (ACADEMIC CALLOUT / NOTE)
% Không dùng box màu mè kiểu brochure, thiết kế thanh lịch với viền trái nhẹ.
% Cú pháp: \academicnote{Nội dung ghi chú}
% ------------------------------------------------------------------------------
\newcommand{\academicnote}[1]{
    \begin{center}
    \begin{minipage}{0.95\linewidth}
        \vspace{4pt}
        \noindent\makebox[\linewidth]{\rule{\linewidth}{0.5pt}}\\[-3pt]
        \small\textbf{Ghi chú:} #1\\[-5pt]
        \noindent\makebox[\linewidth]{\rule{\linewidth}{0.5pt}}
        \vspace{4pt}
    \end{minipage}
    \end{center}
}

% ------------------------------------------------------------------------------
% 4. ĐỊNH NGHĨA KÝ HIỆU TOÁN THƯỜNG DÙNG TRONG HỌC MÁY / AI
% ------------------------------------------------------------------------------
\newcommand{\vect}[1]{\boldsymbol{#1}}       % Ký hiệu vector in đậm (hỗ trợ cả chữ Hy Lạp và Latin)
\newcommand{\mat}[1]{\boldsymbol{#1}}        % Ký hiệu ma trận in đậm
\newcommand{\argmax}{\mathop{\mathrm{argmax}}\limits}
\newcommand{\argmin}{\mathop{\mathrm{argmin}}\limits}
\newcommand{\E}{\mathbb{E}}              % Kỳ vọng toán học
\newcommand{\R}{\mathbb{R}}              % Tập số thực
